\begin{table}[H]
    \centering
    \fontsize{10}{12.5}\selectfont
    \setlength{\tabcolsep}{4pt}
    \caption{Tramos viales dirigidos, distancias métricas y costo de impedancia en el Cusco.}
    \label{tab:aristas_cusco}
    \begin{tabularx}{\textwidth}{@{}ccXrrr@{}}
        \toprule
        \textbf{Origen} & \textbf{Destino} & \textbf{Vía Urbana Asociada} & \textbf{Distancia (m)} & \textbf{Pendiente (\%)} & \textbf{Tiempo Est. (s)} \\
        \midrule
        A & B & Calle Espaderos / Portal Carnicerías & 180 & +1.5 & 44.8 \\
        A & C & Calle Loreto / Pampa del Castillo & 450 & -2.0 & 108.0 \\
        B & D & Calle Garcilaso / Calle Marqués & 220 & +3.0 & 56.8 \\
        B & C & Calle San Bernardo / Almagro & 400 & -1.0 & 96.0 \\
        C & F & Calle Ahuacpinta / Limacpampa Grande & 350 & +0.5 & 36.5 \\
        C & E & Calle Matará / Tres Cruces de Oro & 600 & +2.0 & 151.2 \\
        D & E & Calle Santa Clara / Arco Santa Clara & 250 & +2.5 & 63.8 \\
        D & C & Calle Marqués / Mantas & 500 & -2.5 & 120.0 \\
        E & F & Calle Nueva Baja / Av. Tullumayo & 700 & -1.0 & 72.0 \\
        \bottomrule
    \end{tabularx}
    \vspace{0.15cm}
    \raggedright
    \textit{Nota.} Tiempo de viaje estimado según la función de costo e impedancia por pendiente andina ($\alpha = 2.5$).
\end{table}
